\section{Background}
\textbf{Transformer-based LLM Inference.}
Decoder-only LLMs are composed of stacked Transformer blocks, each consisting primarily of a self-attention module and an MLP. LLM inference consists of a \emph{prefill} phase followed by a \emph{decode} phase. During prefill, the prompt is processed in parallel and the resulting keys and values at each layer are stored in a KV cache. During decode, the model generates one token at a time; for each new token, attention must read the cached KV of the full prior context of that request. In long-context inference, this repeated KV-cache access dominates attention cost. This creates an important asymmetry in batched serving: MLP layers are dominated by shared model-weight accesses and typically benefit from batching, whereas attention remains tied to the per-request cached context and is therefore more memory-bound. Accordingly, decoding may use either \emph{full-context attention}, which attends to the entire cached context, or \emph{sparse-context attention}, which attends only to a selected sparse subset of that context. For request \(i\), we denote this selected sparse context by \(\mathcal{S}_i\). Sparse-context attention reduces attention cost, but may degrade accuracy when \(\mathcal{S}_i\) omits context positions that are important for the current token.

\textbf{Sparse-Context Self-Speculative Decoding.}
Self-speculative decoding alternates between \emph{drafting} and \emph{verification}. During drafting, request \(i\) generates speculative tokens using sparse-context attention over \(\mathcal{S}_i\). During verification, the same target model runs a full-context forward pass over the speculative prefix and compares the drafted tokens against the corresponding full-context outputs. Let \(d_i\) denote the current \emph{draft length}, i.e., the number of buffered drafted tokens that have not yet been verified. If request \(i\) verifies at draft length \(d_i\), and \(a_i \in \{0,\dots,d_i\}\) drafted tokens are accepted, then the accepted drafted prefix, together with the verified next token, is appended to the request's context. The request then resets its draft length and begins accumulating a new speculative prefix. Existing batched speculative decoding systems typically use synchronized execution: requests in the batch share global draft and verify phases under a common speculation schedule. All requests draft together until a batch-level condition triggers verification, and all requests then verify together. This design simplifies execution, but it forces heterogeneous requests to advance in lockstep even when their optimal verification depths differ.
